% !TeX root = main.tex
%% Body of the Mathlib-containment rubric and judge prompt (no heading: the
%% caller supplies one).
\label{app:containment-prompt}

The five-level rubric is:

\begin{lstlisting}[basicstyle=\ttfamily\scriptsize,frame=single,breaklines=true]
5 -- Fully contained. The exact statement (or a direct, trivially equivalent reformulation) exists in Mathlib as a named declaration. A user could cite it directly, or with a thin wrapper.

4 -- Substantially contained. The statement is provable in a few lines from existing Mathlib lemmas, or a strictly more general version is in Mathlib and the target is an immediate specialization.

3 -- Partially contained. Core ingredients (definitions and key supporting lemmas) exist in Mathlib, but the headline statement itself is not there and would require nontrivial assembly.

2 -- Minimally contained. Only background definitions or unrelated prerequisites are available; the substantive content is missing.

1 -- Not contained. Mathlib does not contain the statement or the necessary specialized definitions.

Operational guidance:
- Judge library containment, not mathematical validity, beauty, novelty, importance, or apparent difficulty.
- "A few lines" means a short Lean proof assembled directly from existing declarations, not merely a short informal mathematical argument.
- Use score 5 only when there is a concrete matching named declaration (possibly with symmetry, simplification, notation changes, or an immediate specialization). Identify that declaration whenever possible.
- A verified proof and its referenced Mathlib declarations are supplied as evidence. They are not exhaustive: a generated proof may overlook a more direct existing theorem, and proof length is not by itself dispositive.
- Do not infer a high score merely because the theorem is elementary. Conversely, do not infer a low score merely because its generated proof is long.
- If the exact declaration name is uncertain, report that uncertainty and avoid score 5 unless the API match is otherwise unmistakable.
\end{lstlisting}

Each item is judged with:

\begin{lstlisting}[basicstyle=\ttfamily\scriptsize,frame=single,breaklines=true]
You are an impartial expert judge of Lean 4 Mathlib library containment.

Apply the supplied rubric exactly. Estimate how much of the target theorem is already contained in Mathlib. The theorem has been verified, so do not judge whether it is true. Do not judge its mathematical novelty, importance, elegance, or difficulty.

The proof body and referenced declarations are supporting evidence, not an exhaustive Mathlib search. A generated proof can miss a direct existing theorem. Treat score 5 as requiring a concrete matching named declaration or an unmistakably direct API match. Distinguish a genuinely few-line Mathlib wrapper (score 4) from nontrivial assembly of ingredients (score 3).

<rubric>
[RUBRIC ABOVE]
</rubric>

<target_theorem>
[STATEMENT]
</target_theorem>

<verified_proof_body lines="[LINE COUNT]">
[PROOF BODY]
</verified_proof_body>

<referenced_mathlib_declarations>
[NAMES AND TYPES, OR NO RECOVERED LIST]
</referenced_mathlib_declarations>

Return exactly one JSON object and no Markdown or surrounding commentary:
{
  "score": <integer 1-5>,
  "matching_declaration": "<concrete Mathlib declaration for score 5, best candidate otherwise, or none>",
  "supporting_declarations": ["<zero or more important declaration names>"],
  "rationale": "<concise explanation tied to the rubric>",
  "confidence": "<high | medium | low>"
}
\end{lstlisting}
